\section{Consequences and limits of the characterization}\label{sec:discussion}

The selection formula recovers two immediate upper bounds by choosing
the retained set to be all or none of the magnitude space:
\[
 R_{\mu,\rho}(B)\lesssim
 \min\left\{\rho,\sqrt{\rho\Gamma_\mu/B},
 \int\min\{1,(d\sqrt B)^{-1}\}\,\mu(dd)\right\}.
\]
The full formula additionally allows different regions of the same law
to use different treatments. Its lower proof explains why contributions
from several irregular components may need to be combined.

For a single magnitude $\mu=\delta_d$, hard assignments leave two choices
and give
\[
 R_{\delta_d,\rho}(B)\asymp
       \min\{\rho,\sqrt{\rho/(Bd^2)}\}.
\]
For a general law, the finite-prefix selection values $k_J$ in
\eqref{eq:canonical-selection-prefix} have an explicit tail-mass error
bound \eqref{eq:canonical-prefix-limit}, and their limit lies between
$\mathcal K/2$ and $\mathcal K$. This is a finite approximation using only
known band masses. Exhaustive selection is possible, but an efficient
general selection algorithm has not been established here.

The scalar characterization concerns the full minimax risk, whereas
$P_N$ bounds worst bias and worst variance separately. Accordingly, its
exact dual is an identity for a certificate program, and the statistical
connection is a comparison inequality. The convergent genie brackets
can certify that a chosen cap suffices for a particular input without
using the large universal cap. They do not give a uniform runtime bound.

The remaining quantitative questions include sharp comparison constants,
efficient selection for structured laws, and rates of cap convergence.
The present theorem assumes known $\mu$, Bernoulli observations, and a
uniform expected observation budget. Extensions to an unknown magnitude
law, other observation kernels, or almost-sure budgets require separate
arguments.

This consolidated manuscript includes the specialized proof chain rather
than relying on the earlier working notes. Separate model audits found
no remaining substantive inconsistency in the reviewed chain; exact
arithmetic verifies the finite numerical certificates. These checks do
not constitute external refereeing or proof-assistant verification of the
general theorems. The related-work comparison is targeted and does not
certify publication priority.


\subsection{Two scope examples}

\paragraph{A noninterval optimal assignment.}
Let $\varepsilon=1/1024$, put masses $(1-\varepsilon)/2$, $\varepsilon$,
and $(1-\varepsilon)/2$ at $1/2$, $1/4$, and $1/8$, respectively, and
take $\rho=1/4$, $h=B^{-1/2}=1/1024$.
All three atoms lie above $h$. Enumerating the eight hard retained sets
shows that the unique minimum retains $\{1/2,1/8\}$ and treats $1/4$
ordinarily. The following rational enclosures for the inner cost divided
by $h$ make the strict comparison explicit:
\begin{center}
\begin{tabular}{@{}lc@{}}
\toprule
Retained atoms & Enclosure for cost divided by $h$\\
\midrule
$\varnothing$ & $(4.99902,4.99903)$\\
$\{1/2\}$ & $(4.70676,4.70677)$\\
$\{1/4\}$ & $(5.05761,5.05762)$\\
$\{1/2,1/4\}$ & $(4.72155,4.72156)$\\
$\{1/8\}$ & $(3.82997,3.82998)$\\
$\{1/2,1/8\}$ & $(3.22748,3.22749)$\\
$\{1/4,1/8\}$ & $(3.83271,3.83272)$\\
$\{1/2,1/4,1/8\}$ & $(3.22943,3.22944)$\\
\bottomrule
\end{tabular}
\end{center}
These follow by substituting the rational weights into $\mathcal G$.
For example, after reducing each logarithm argument to $1\le x\le2$,
use $y=(x-1)/(x+1)$ and
\[
 0\le\log x-2\sum_{k=0}^{m-1}\frac{y^{2k+1}}{2k+1}
 \le\frac{2y^{2m+1}}{(2m+1)(1-y^2)},
\]
together with rational bounds for $\log2$ from the same series and
square-root bounds checked by squaring. Already $m=90$ gives intervals
much narrower than those displayed. The companion verifier
\texttt{checks/verify\_revisions.py} performs these exact calculations.
Since the outer cap is inactive, the original functional has the same
optimizer. No interval can contain both retained atoms while excluding
the intervening positive-mass atom.

\paragraph{Infinite inverse moments and a slow tail.}
Let $\mu\{2^{-j}\}=1/[j(j+1)]$ for $j\ge1$. Both inverse moments
of orders one and two are infinite. At $B=4^J$, ignore $j>J$ and
treat the first $J$ atoms ordinarily. This gives
\[
 \mathcal K(B)\le\frac1{J+1}
       +\sum_{j=1}^J\frac{2^{j-J}}{j(j+1)}
       =\frac1{J+1}+O(J^{-2}).
\]
For fixed $\rho$ and sufficiently large $J$, the ignored-tail allocation
gives $\Psi(2^{-J})\ge1/(J+1)$. Thus the upper and lower scales match
even though the full-law inverse moments diverge. This example uses
only the vanishing tail probability, not an integrable full-law cost.
